\documentclass[a4paper,12pt]{article}
\usepackage[T2A]{fontenc}
\usepackage[utf8]{inputenc}
\usepackage[russian]{babel}
\usepackage{amsmath,amssymb,mathtools}
\usepackage{hyperref}
\usepackage{enumitem}
\usepackage{booktabs}
\usepackage{geometry}
\geometry{top=2cm, bottom=2cm, left=2.5cm, right=2.5cm}

\title{Конспект лекции 2: \\ Свёрточные нейронные сети}
\author{Липкин С.М.}
\date{}

\begin{document}
\maketitle

\section{Мотивация}

Полносвязные (FC) слои плохо подходят для обработки изображений по нескольким причинам:
\begin{itemize}
    \item Огромное число параметров: изображение $224 \times 224 \times 3$ даёт 150\,528 входных нейронов. Один скрытый слой на 1024 нейрона — 150 миллионов параметров.
    \item Нет инвариантности к сдвигу: кошка в левом углу и кошка в правом углу — разные активации.
    \item Нет локальности: каждый пиксель влияет на каждый нейрон.
\end{itemize}

Свёрточные нейронные сети (CNN) решают эти проблемы через:
\begin{itemize}
    \item \textbf{Локальную связность} (local connectivity): каждый нейрон видит только локальную область изображения (receptive field).
    \item \textbf{Разделение весов} (weight sharing): один и тот же фильтр применяется ко всем участкам. Это даёт инвариантность к сдвигу и в 10--100 раз меньше параметров.
    \item \textbf{Иерархию признаков}: первые слои учат края и текстуры, следующие — формы и части объектов.
\end{itemize}

\section{Что такое свёртка}

\subsection{Математическое определение}

Дискретная 1D свёртка:
\[
(f * g)[n] = \sum_{m=-\infty}^{\infty} f[m] \cdot g[n-m]
\]

На практике в CNN используется \textbf{кросс-корреляция} (ядро не переворачивается):
\[
S[i,j] = \sum_{u=0}^{K_h-1} \sum_{v=0}^{K_w-1} I[i+u, j+v] \cdot K[u, v]
\]

где:
\begin{itemize}
    \item $I$ — входной сигнал (изображение, feature map) размера $H_{\text{in}} \times W_{\text{in}}$
    \item $K$ — ядро (kernel, filter) размера $K_h \times K_w$
    \item $S$ — выходной сигнал (feature map)
\end{itemize}

\subsection{Интуиция}

Ядро «скользит» по изображению с шагом $S$ (stride). В каждой позиции вычисляется скалярное произведение между ядром и локальной областью изображения. Результат — карта активаций (feature map), показывающая, \textbf{где} ядро «нашло» искомый паттерн.

Если ядро — детектор горизонтальных краев, то feature map будет иметь высокие значения там, где есть горизонтальные границы.

\section{Свёрточная арифметика}

\subsection{Основная формула}

Размер выхода свёрточного слоя:
\[
O = \left\lfloor \frac{I + 2P - K}{S} \right\rfloor + 1
\]

где:
\begin{itemize}
    \item $I$ — размер входа (по одному измерению)
    \item $K$ — размер ядра
    \item $P$ — padding (количество нулевых пикселей по краям)
    \item $S$ — stride (шаг)
    \item $O$ — размер выхода
\end{itemize}

Если по разным осям padding или stride разный, формула применяется отдельно к каждой оси:
\[
H_{\text{out}} = \left\lfloor \frac{H_{\text{in}} + 2P_h - K_h}{S_h} \right\rfloor + 1
\]
\[
W_{\text{out}} = \left\lfloor \frac{W_{\text{in}} + 2P_w - K_w}{S_w} \right\rfloor + 1
\]

\subsection{Типы padding}

\begin{itemize}
    \item \textbf{Valid} ($P=0$): $O = (I - K) / S + 1$ — выход сжимается.
    \item \textbf{Same} ($P = \lfloor K/2 \rfloor$): $O = I / S$ — размер сохраняется (при $S=1$).
    \item \textbf{Full} ($P = K - 1$): $O = I + K - 1$ — вход и ядро полностью перекрываются хотя бы на одном пикселе.
\end{itemize}

\subsection{Примеры}

\begin{itemize}
    \item $I=32$, $K=5$, $S=1$, Valid: $O = 32 - 5 + 1 = 28$
    \item $I=32$, $K=5$, $S=1$, Same: $O = 32$ (при $P=2$)
    \item $I=32$, $K=3$, $S=2$, Valid: $O = \lfloor (32 - 3) / 2 \rfloor + 1 = 15$
\end{itemize}

\subsection{Многоканальная свёртка}

Вход: $H \times W \times C_{\text{in}}$ (например, RGB: 3 канала).

Ядро: $K_h \times K_w \times C_{\text{in}} \times C_{\text{out}}$. Каждый выходной канал — сумма свёрток по всем входным каналам:
\[
S_j = \sum_{i=1}^{C_{\text{in}}} I_i * K_{ij}
\]

Количество параметров слоя:
\[
\text{params} = K_h \cdot K_w \cdot C_{\text{in}} \cdot C_{\text{out}} + C_{\text{out}}
\]

\subsection{Receptive Field}

Размер области входа, «видимой» одним нейроном:
\[
RF_l = RF_{l-1} + (K_l - 1) \cdot \prod_{i=1}^{l-1} S_i
\]

Начальное значение: $RF_0 = 1$.

\section{Базовые блоки CNN}

\subsection{Pooling (субдискретизация)}

\textbf{Max Pooling:}
\[
y_{ij} = \max_{u=1..K, v=1..K} x_{iS+u, jS+v}
\]

Выделяет наиболее активный признак, инвариантен к малым сдвигам. Нет обучаемых параметров.

\textbf{Average Pooling:}
\[
y_{ij} = \frac{1}{K^2} \sum_{u=1..K, v=1..K} x_{iS+u, jS+v}
\]

Более «гладкий» выход.

\textbf{Global Pooling:} pooling по всему пространству ($K = H = W$) $\to$ $1 \times 1$.

Зачем: уменьшение размерности $\to$ меньше параметров $\to$ bigger receptive field без роста вычислений.

\subsection{Транспонированная свёртка (Deconvolution)}

Операция увеличения пространственного разрешения. Если обычная свёртка — от многих к одному, то транспонированная — от одного ко многим.

Формула:
\[
O = (I - 1) \cdot S + K - 2P
\]

\textbf{Важно:} это \textbf{не} обратная свёртка. Это fractionally-strided convolution. Используется в генеративных моделях (GAN, VAE), сегментации (U-Net) и super-resolution.

\subsection{Dilated (Atrous) Convolution}

Увеличивает receptive field без увеличения числа параметров:
\[
K_{\text{effective}} = K + (K - 1) \cdot (d - 1)
\]

где $d$ — dilation rate. Для $3 \times 3$ ядра:
\begin{itemize}
    \item $d = 1$: $K_{\text{eff}} = 3$, RF = $3 \times 3$
    \item $d = 2$: $K_{\text{eff}} = 5$, RF = $5 \times 5$
    \item $d = 4$: $K_{\text{eff}} = 9$, RF = $9 \times 9$
\end{itemize}

Параметров не прибавляется. Используется в DeepLab для семантической сегментации.

\subsection{Depthwise Separable Convolution}

Разделяет стандартную свёртку на два этапа:
\begin{enumerate}
    \item \textbf{Depthwise:} $K \times K$ свёртка для каждого канала отдельно: $\text{params}_{\text{dw}} = K^2 \cdot C_{\text{in}}$
    \item \textbf{Pointwise ($1 \times 1$):} линейная комбинация каналов: $\text{params}_{\text{pw}} = C_{\text{in}} \cdot C_{\text{out}}$
\end{enumerate}

\textbf{Сравнение:} для $3\times3$, $C_{\text{in}}=64$, $C_{\text{out}}=128$:
\begin{itemize}
    \item Обычная: $3^2 \cdot 64 \cdot 128 = 73\,728$ параметров
    \item Depthwise Separable: $3^2 \cdot 64 + 64 \cdot 128 = 8\,768$ параметров
    \item Ускорение: \textbf{8.4x}
\end{itemize}

Используется в MobileNet, Xception, EfficientNet.

\subsection{$1 \times 1$ Свёртка (Pointwise Convolution)}

Линейная комбинация каналов без пространственного взаимодействия:
\[
y_{ij}^{(k)} = \sum_{c=1}^{C_{\text{in}}} w_c^{(k)} \cdot x_{ij}^{(c)}
\]

\textbf{Применения:} изменение числа каналов (bottleneck), снижение размерности перед дорогими $3\times3$ свёртками.

\textbf{Bottleneck в ResNet:} $1\times1$ (256 $\to$ 64) $\to$ $3\times3$ (64 $\to$ 64) $\to$ $1\times1$ (64 $\to$ 256). $3\times3$ работает в 4 раза быстрее (64 канала вместо 256).

\section{Архитектуры CNN}

\subsection{LeNet-5 (LeCun, 1998)}

Первая успешная CNN. Для распознавания рукописных цифр MNIST.

\textbf{Архитектура:}
\begin{itemize}
    \item Вход: $32 \times 32$ grayscale
    \item Conv $5\times5$, 6 каналов $\to$ AvgPool $2\times2$
    \item Conv $5\times5$, 16 каналов $\to$ AvgPool $2\times2$
    \item Conv $5\times5$, 120 каналов (фактически FC)
    \item FC 84 $\to$ FC 10 (softmax)
\end{itemize}

Всего 60K параметров. Паттерн «свёртка $\to$ пулинг $\to$ свёртка $\to$ пулинг $\to$ FC» стал шаблоном для всех последующих CNN.

\subsection{AlexNet (Krizhevsky, 2012)}

Победитель ImageNet 2012. Переломный момент для DL.

\textbf{Ключевые инновации:}
\begin{itemize}
    \item \textbf{ReLU} вместо tanh — решил проблему исчезающего градиента.
    \item \textbf{Data augmentation} — сдвиги, отражения, изменение цвета.
    \item \textbf{Dropout} (0.5 на FC-слоях) — борьба с overfitting.
    \item Local Response Normalization (устарело, заменено BatchNorm).
    \item Обучение на 2 GPU параллельно.
\end{itemize}

\textbf{Архитектура:} 5 свёрточных + 3 FC слоя. Conv: $11\times11$ (96 каналов), $5\times5$ (256), $3\times3$ (384, 384, 256). MaxPool $3\times3$ после первых двух и последней свёртки. FC: 4096 $\to$ 4096 $\to$ 1000.

62M параметров. Время обучения: 5--6 дней. Запустил эру глубокого обучения.

\subsection{VGG (Simonyan \& Zisserman, 2014)}

\textbf{Главная идея:} стек из $3\times3$ свёрток. Два $3\times3$ слоя $\to$ RF $5\times5$. Три $3\times3$ слоя $\to$ RF $7\times7$.

\textbf{Почему это выгоднее одного большого ядра:}
\begin{itemize}
    \item $3\times3$: 9 параметров на канал.
    \item $5\times5$: 25 параметров на канал ($2.78\times$ больше).
    \item $7\times7$: 49 параметров на канал ($5.44\times$ больше).
    \item Дополнительные нелинейности между слоями.
\end{itemize}

\textbf{VGG16:} Conv $3\times3$: 64 $\to$ 128 $\to$ 256 $\to$ 512 $\to$ 512, MaxPool $2\times2$ после каждой группы, 3 FC слоя. 138M параметров. Показал, что глубина важнее ширины.

\subsection{Inception / GoogLeNet (Szegedy, 2014)}

\textbf{Проблема:} какой размер ядра выбрать?

\textbf{Решение Inception module:} используем \textbf{все сразу} — параллельные ветки с $1\times1$, $3\times3$, $5\times5$ свёртками и $3\times3$ MaxPool, конкатенация по каналам.

$1\times1$ свёртки перед дорогими $3\times3$ и $5\times5$ играют роль bottleneck — снижают размерность.

\textbf{Результат:} 6.8M параметров (VGG: 138M) при сопоставимом качестве. Использовал auxiliary classifiers для борьбы с исчезающим градиентом.

\subsection{ResNet (He et al., 2015)}

\textbf{Проблема:} при $> 20$ слоях качество падает (degradation problem). Не из-за overfitting — и train, и test loss растут.

\textbf{Решение: Residual Learning}
\[
H(x) = F(x) + x
\]

Слой учится \textbf{добавке} (residual) $F(x)$ к тождественному отображению $x$. Если слой не нужен: $F(x) \to 0$, $H(x) \to x$. Градиент течёт через shortcut напрямую.

\textbf{Влияние на градиент:}
\[
\frac{\partial \mathcal{L}}{\partial x_l} = \frac{\partial \mathcal{L}}{\partial x_{l+1}} \left(1 + \frac{\partial F(x_l)}{\partial x_l}\right)
\]

Единичная матрица сохраняет градиент — vanishing gradient практически устранён.

\textbf{Bottleneck block:} $1\times1$ (256 $\to$ 64) $\to$ $3\times3$ (64 $\to$ 64) $\to$ $1\times1$ (64 $\to$ 256).

\textbf{Результат:} 152 слоя, 3.57\% top-5 error на ImageNet. «Рабочая лошадка» до сих пор.

\subsection{DenseNet (Huang et al., 2017)}

\textbf{Идея:} каждый слой получает на вход feature maps со \textbf{всех предыдущих слоёв}:
\[
x_l = H_l([x_0, x_1, \dots, x_{l-1}])
\]

\textbf{Преимущества:}
\begin{itemize}
    \item Максимальный поток информации и градиентов — каждый слой имеет прямой доступ к градиенту лосса.
    \item Естественная регуляризация — нет дублирующих признаков.
    \item Меньше параметров, чем у ResNet при схожем качестве (DenseNet-121: 8M vs ResNet-50: 25M).
\end{itemize}

Недостаток: большой расход памяти (нужно хранить все активации для backpropagation).

\subsection{MobileNet (2017--2018)}

Оптимизированы для мобильных устройств:
\begin{itemize}
    \item Depthwise Separable Convolution везде.
    \item Width multiplier $\alpha$ — пропорциональное уменьшение числа каналов.
    \item Resolution multiplier $\rho$ — уменьшение разрешения входа.
\end{itemize}

\textbf{MobileNetV2:} Inverted Residual (узкий $\to$ широкий $\to$ узкий) + Linear bottleneck.

\subsection{EfficientNet (Tan \& Le, 2019)}

\textbf{Ключевая идея:} композитное масштабирование. Обычно масштабируют \textbf{один} параметр (глубину, ширину или разрешение). EfficientNet масштабирует \textbf{все три} сбалансированно:
\[
d = \alpha^\phi, \quad w = \beta^\phi, \quad r = \gamma^\phi
\]
где $\alpha\beta^2\gamma^2 \approx 2$ (FLOPs растут в $2^\phi$ раз).

EfficientNet-B7: 66M параметров, достигает SOTA при \textbf{в 10x меньше FLOPs}, чем предыдущие SOTA.

\section{Сводка архитектур}

\begin{center}
\begin{tabular}{lrrrl}
\toprule
\textbf{Архитектура} & \textbf{Год} & \textbf{Парам.} & \textbf{Top-5 err.} & \textbf{Ключевая идея} \\
\midrule
LeNet-5 & 1998 & 60K & — & Первая CNN \\
AlexNet & 2012 & 62M & 16.4\% & ReLU, Dropout, GPU \\
VGG16 & 2014 & 138M & 8.8\% & $3\times3$ стеки \\
GoogLeNet & 2014 & 6.8M & 6.7\% & Inception modules \\
ResNet-152 & 2015 & 60M & 3.6\% & Skip connections \\
DenseNet-121 & 2017 & 8M & 5.6\% & Dense connections \\
MobileNetV2 & 2018 & 3.5M & — & Depthwise sep. \\
EfficientNet-B7 & 2019 & 66M & 2.2\% & Compound scaling \\
\bottomrule
\end{tabular}
\end{center}

\section{Заключение}

Эволюция архитектур CNN: LeNet задал паттерн, AlexNet показал масштаб, VGG — глубину, Inception — ширину, ResNet решил degradation problem, DenseNet максимизировал поток информации, MobileNet оптимизировал для устройств, EfficientNet нашёл баланс масштабирования.

Несмотря на доминирование Transformer в NLP и Vision (ViT), CNN остаются основой для множества приложений: детекция, сегментация, генеративные модели, обработка изображений на устройствах.

\end{document}
